\chapter{Diseño experimental, selección y cartera}
\label{chap:diseno-experimental}

El diseño separa de forma estricta el aprendizaje predictivo de la traducción económica. La
configuración ganadora sólo se elige con Rank-IC robusto, calculado sobre 2015--2024; costes,
turnover, alfa, perfiles y rendimiento frente a SPY se calculan después como diagnósticos. Esta
separación evita seleccionar retrospectivamente una cartera porque su rentabilidad conocida fue
alta, una práctica especialmente problemática cuando se prueban muchas alternativas
(Harvey et al., 2016; López de Prado et al., 2014).

\section{Optimización secuencial y regla de decisión}

El catálogo cerrado se recorre en cuatro fases: temporal, representación, modelo y meta. Se parte
de un \textit{baseline}; en cada fase se modifica una variable y se comparan sus candidatos con el
incumbente acumulado. Así, el coste científico es una suma de alternativas, no una búsqueda
cartesiana opaca. Las 71 configuraciones ejecutadas en la pasada de referencia quedan registradas
junto con cada elección, y las de las dos pasadas anteriores conservan sus propios registros.

Un candidato debe tener observaciones suficientes, no caer por debajo de $-0,02$ en ninguna era
disponible y superar una puerta pareada: dominar al incumbente (ventaja media positiva y mejor en
más de la mitad de cohortes) o ser no inferior con límite inferior bootstrap al 90\% superior a
$-0,01$. Las diferencias se calculan por periodo mensual, lo que permite comparar lags distintos
sin confundir el desplazamiento de sus fechas de ejecución con ausencia de evidencia.

\begin{figure}[H]
  \centering
  \begin{tikzpicture}[node distance=5mm, >=Latex, every node/.style={font=\small,align=center}]
    \node[draw,rounded corners,fill=blue!8,minimum width=27mm,minimum height=10mm] (t) {Temporal};
    \node[draw,rounded corners,fill=teal!10,right=of t,minimum width=30mm,minimum height=10mm] (r) {Representación};
    \node[draw,rounded corners,fill=orange!12,right=of r,minimum width=25mm,minimum height=10mm] (m) {Modelo};
    \node[draw,rounded corners,fill=green!10,right=of m,minimum width=25mm,minimum height=10mm] (x) {Meta};
    \node[draw,rounded corners,fill=gray!15,below=8mm of r,minimum width=48mm,minimum height=10mm] (w) {Ganador congelado};
    \node[draw,rounded corners,fill=yellow!18,right=of w,minimum width=38mm,minimum height=10mm] (d) {Robustez y cartera\\diagnósticas};
    \draw[->,thick] (t)--(r); \draw[->,thick] (r)--(m); \draw[->,thick] (m)--(x);
    \draw[->,thick] (x.south)--++(0,-5mm)-|(w.north);
    \draw[->,thick] (w)--(d);
  \end{tikzpicture}
  \caption{Protocolo de selección. Sólo las cuatro fases superiores pueden modificar el ganador;
  la evidencia económica se calcula una vez congelada la señal.}
  \label{fig:protocolo}
\end{figure}

La Tabla~\ref{tab:decisiones} hace auditable esa traza para la pasada de referencia. Conviene
explicar qué columna se tabula y por qué, porque la elección no es inocente. El Rank-IC acumulado
tras cada paso --- la representación habitual, en forma de escalera --- sería en esta pasada una
columna constante: 0,1005 repetido dieciséis veces. No por un defecto de la medición, sino porque la
tercera pasada parte del ganador de la segunda y casi todas las variables llegan ya en su mejor
valor conocido, de modo que la comparación pareada no encuentra motivo para moverlas.

Lo informativo, entonces, no es el valor que se conserva sino \textbf{la alternativa que se
descarta}: cuánto Rank-IC habría costado adoptarla. Esa es la columna que se tabula, ordenada de
mayor a menor coste.

\begin{table}[H]
  \centering
  \small
  \caption{Qué alternativa se rechazó en cada decisión y cuánto habría costado adoptarla.}
  \label{tab:decisiones}
  \begin{tabular}{lllrrl}
\toprule
Variable & Ganador & Mejor alternativa & Su Rank-IC & Coste & Regla \\
\midrule
neutralize\_by\_sector & False & True & 0,0576 & -0,0429 & Rank-IC robusto \\
recency\_weighting & off & linear & 0,0632 & -0,0373 & Rank-IC robusto \\
target\_horizon\_months & 12 & 6 & 0,0684 & -0,0322 & Rank-IC robusto \\
feature\_preset & all & core & 0,0716 & -0,0290 & Rank-IC robusto \\
snapshot\_step\_months & 1 & 3 & 0,0916 & -0,0089 & Rank-IC robusto \\
meta\_method & stacked rolling free & stacked rolling bounded & 0,1005 & -0,0085 & Rank-IC robusto \\
train\_lookback\_years & 8 & 12 & 0,0953 & -0,0052 & Rank-IC robusto \\
feature\_weighting\_mode & oos stability prune & model native & 0,0959 & -0,0047 & Rank-IC robusto \\
fundamental\_momentum & True & False & 0,0967 & -0,0038 & Rank-IC robusto \\
winsorization & 0.0 & 0.01 & 0,0984 & -0,0021 & Rank-IC robusto \\
lgbm\_max\_depth & 3 & 4 & 0,0992 & -0,0014 & Rank-IC robusto \\
lgbm\_learning\_rate & 0.03 & 0.05 & 0,0995 & -0,0010 & Rank-IC robusto \\
lgbm\_n\_estimators & 100 & 200 & 0,0997 & -0,0008 & Rank-IC robusto \\
execution\_lag\_days & 60 & 30 & 0,1000 & -0,0005 & Empate: simplicidad \\
max\_features\_per\_agent & 20 & 12 & 0,1000 & -0,0005 & Rank-IC robusto \\
market\_regime\_feature & False & True & 0,1001 & -0,0004 & Rank-IC robusto \\
lgbm\_min\_child\_samples & 20 & 50 & 0,1009 & 0,0004 & Empate: simplicidad \\
\bottomrule
\end{tabular}

  \caption*{Fuente: \artefacto{decisions.json}. El coste es el Rank-IC de la mejor alternativa menos
  el del ganador; un coste \emph{positivo} indica que la alternativa medía mejor y que el ganador se
  decidió por otra vía, lo que sólo ocurre en los empates técnicos.}
\end{table}

Leída así, la tabla separa con claridad tres grupos. Cuatro decisiones eran caras de equivocar:
neutralizar por sector habría costado $-0{,}0429$ de Rank-IC, aplicar pesos de recencia $-0{,}0373$,
acortar el horizonte a seis meses $-0{,}0322$ y reducir el conjunto de variables a \texttt{core}
$-0{,}0290$. Son las decisiones estructurales, y las cuatro se resolvieron con holgura.

Un segundo grupo, el más numeroso, tiene costes por debajo de 0,005: los hiperparámetros de LightGBM
(profundidad, número de árboles, tasa de aprendizaje), la winsorización o el número máximo de
variables por agente. Ahí la evidencia apenas distingue entre alternativas, y conviene decirlo
abiertamente: el sistema no debe su capacidad predictiva al ajuste fino del modelo.

El tercer grupo son las dos decisiones resueltas por empate técnico, que la última columna marca de
forma explícita. En \artefacto{execution\_lag\_days} la alternativa de 30 días quedó $0{,}0005$ por
debajo del ganador y en \artefacto{lgbm\_min\_child\_samples} el valor 50 medía $0{,}0004$
\emph{mejor}. Es decir, en el segundo caso el ganador no es el valor que mejor midió. Ninguna de las
dos debe presentarse como hallazgo empírico, y por eso figuran con su regla al lado.

\subsection*{Por qué la escalera es plana}

El resultado de la tercera pasada es contundente y va en contra de la intuición habitual sobre este
tipo de búsquedas: \textbf{sólo una de las diecisiete decisiones cambió el incumbente}. Las
dieciséis restantes confirmaron el valor de partida, con una ventaja pareada exactamente nula. El
Rank-IC permanece en 0,1005 durante toda la pasada y sólo se mueve en el último paso,
\artefacto{meta\_method}, que lo lleva a 0,1090.

Esa planitud no es un defecto del procedimiento sino la consecuencia esperable de encadenar
estudios. La tercera pasada parte del ganador de la segunda, de modo que casi todas las variables
llegan ya en su mejor valor conocido y la comparación pareada no encuentra motivo para moverlas.
Una escalera plana en la última pasada de una cadena significa que la cadena ha convergido; la
primera pasada, que partía del baseline recomendado del catálogo, presentaba un perfil muy distinto,
como documenta la Sección~\ref{sec:cadena-studies}.

Este patrón es una defensa parcial frente a la objeción de sobreajuste por multiplicidad. Un
protocolo que aceptase cualquier mejora aparente habría desplazado al incumbente muchas más veces,
persiguiendo diferencias dentro del ruido. Que dieciséis de diecisiete decisiones terminen en «no
cambiar» indica que las puertas de elegibilidad están operando. No elimina el problema --- el
Deflated Sharpe del Capítulo~\ref{chap:resultados} sigue penalizando por las configuraciones
ensayadas, y la cadena eleva ese recuento en lugar de reducirlo ---, pero acota su alcance: la
multiplicidad afecta sobre todo a la afirmación económica, no a la traza de decisiones predictivas.

La regla de simplicidad que resolvió esas dos decisiones opera bajo un umbral de empate técnico de
0,002 en ventaja pareada. Es una preferencia declarada de antemano --- ante evidencia equivalente,
se elige la opción más simple o la más conservadora --- y no una medición, lo que delimita
exactamente qué puede afirmarse sobre los valores que produce.

\subsection*{El catálogo cerrado como garantía}

La afirmación de que nada se optimizó fuera del catálogo sólo es verificable si el catálogo es
público. El Anexo~\ref{anexo:catalogo} reproduce las treinta y tres variables con su rejilla
completa, su etapa, si son predictivas y el valor que tomó el ganador. Su función en este documento
no es descriptiva sino probatoria: cualquier lector puede comprobar que el valor ganador de cada
variable pertenece a su rejilla declarada, y que ninguna variable ajena aparece en la configuración
final.

Doce de esas treinta y tres variables son de cartera y están marcadas como no predictivas. La
distinción es la que sostiene el Capítulo~\ref{chap:resultados}: se ejecutan después de congelar el
ganador, se reportan como diagnóstico y no pueden modificarlo. Es lo que convierte el barrido de la
Sección~\ref{sec:cartera-rejilla} en un experimento controlado en lugar de una segunda ronda de
optimización.

\section{Construcción de cartera}

La cartera del ganador contiene 12 posiciones, aplica pesos proporcionales al alfa esperado, una
política de efectivo por oportunidad con tope del 25\%, comisión de 5 pb y \textit{slippage} de 10
pb. SPY se usa sólo como benchmark; nunca es una posición. Los umbrales se expresan en puntos
básicos anuales de alfa esperado para poder compararlos entre horizontes. Una rotación requiere que
la diferencia esperada cubra el coste de ida y vuelta y un margen de 50 pb; una compra nueva usa
histéresis. Estas reglas pretenden preservar señal y controlar rotación, no volver a optimizar el
modelo por rentabilidad.

La era 2025--2026 fue aislada antes de aplicar la selección. Sus observaciones no participan ni en
los pesos del meta ni en la elección de parámetros; su papel es exclusivamente de confirmación
posterior.

\section{Por qué la selección predictiva es secuencial}
\label{sec:seleccion-secuencial}

Una búsqueda cartesiana de todos los valores posibles permitiría encontrar combinaciones muy
ajustadas a la historia, dificultaría explicar de dónde surge el ganador y multiplicaría las
pruebas. El protocolo fija un orden: temporalidad, representación, modelo y meta. Cada fase hereda
el incumbent anterior, de modo que una decisión posterior no puede reabrir silenciosamente una
elección temporal. Si dos alternativas difieren menos de 0,002 en ventaja pareada, se prefiere la
más simple; si tampoco hay evidencia suficiente, se conserva el incumbent.

Conviene precisar el alcance de este argumento, porque el trabajo emplea después una búsqueda
cartesiana y sería contradictorio no explicar por qué. El razonamiento anterior vale para el plano
\textbf{predictivo}, y vale por dos motivos que no son intercambiables. El primero es estadístico:
la selección predictiva se decide sobre el Rank-IC, la misma magnitud que después se reporta como
resultado principal, de modo que cada comparación adicional infla el mejor valor observado y
contamina precisamente la cifra que sostiene el trabajo. El segundo es expositivo: el objetivo del
Model Study no es solo encontrar una configuración buena, sino poder explicar de dónde sale cada
decisión, y una rejilla de miles de puntos devuelve un ganador que nadie puede narrar.

Ninguno de los dos motivos se aplica al plano de la \textbf{cartera}. Las reglas de cartera no
tocan el modelo ni modifican el Rank-IC: operan sobre una señal ya congelada, así que optimizarlas
no puede inflar la métrica predictiva que el documento reporta. Y su elección no necesita narrarse
como una escalera de decisiones porque no lo es: las variables de cartera interactúan entre sí, y un
protocolo secuencial no solo sería innecesario sino directamente inadecuado, como se argumenta en el
Capítulo~\ref{chap:desarrollo-cartera}. Queda un riesgo real, la multiplicidad de comparaciones, que
no se disimula: se declara en el Capítulo~\ref{chap:limitaciones} y se acota por construcción
impidiendo que la rejilla llegue a calcular la era reservada.

\section{Ejemplo de comparación pareada}

Supóngase que un candidato y el incumbent producen Rank-IC sobre las mismas 40 cohortes. En cada
fecha se resta el Rank-IC del incumbente al del candidato y se re-muestrean bloques de doce cohortes
para preservar la dependencia inducida por el horizonte anual. Si la ventaja media es 0,0208, el
candidato mejora el 57,5\% de cohortes y el IC bootstrap al 90\% es $[0,0108;0,0372]$, hay evidencia
de dominancia. Éste es el caso real de la cadencia mensual frente al baseline trimestral, tal como
se decidió en la primera pasada de la cadena.

La rentabilidad de cartera no influye en este paso. Un CAGR alto puede surgir de concentración,
exposición accidental a una acción extrema o efectivo, sin demostrar que se ordene mejor el
universo. El horizonte de seis meses ilustra un rechazo: su ventaja pareada fue $-0,0265$ y el
intervalo $[-0,0552;-0,0109]$. Mostrar este resultado evita presentar la selección como una
confirmación automática de la configuración preferida.

\section{De rango a alfa esperado}

La cartera necesita una magnitud económica, no sólo un percentil. Sobre cohortes cerradas se estima
una relación entre ventiles de \texttt{meta\_rank} y retorno excedente anualizado. Se usan veinte
grupos para no ajustar una curva a medias de muy pocas acciones, aunque la función se evalúa sobre
el rango continuo. Si la pendiente no es creciente, la evidencia se amplía en cascada a dieciséis
trimestres y al histórico completo; la salvaguarda final queda registrada como supuesto y no se
presenta como estimación empírica.

\section{Ejemplo de decisión y coste}

Con horizonte de doce meses, comisión de 5 pb y slippage de 10 pb, una rotación completa cuesta 30
pb. Si la peor posición tiene alfa esperado 60 pb, una candidata externa 250 pb y el margen de
rotación es 50 pb, la mejora de 190 pb supera 80 pb y se rota. Si la candidata ofrece 120 pb, la
mejora no cubre coste y margen, por lo que se conserva la posición. El ranking ordena preferencias;
la operación necesita además una ventaja neta.

Las entradas nuevas incorporan histéresis y el efectivo rinde 0\%. Bajo
\texttt{opportunity\_cash}, una plaza puede permanecer en efectivo sólo si respeta el tope de 25\%
y el suelo de diversificación. Un mínimo de tenencia, si se activa, bloquea ventas por rotación o
caída de alfa hasta cumplir la fracción de horizonte configurada. Estas reglas controlan fricción y
estabilidad; no pueden alterar el ganador predictivo.
